\documentclass[12pt]{article}
\usepackage{amsmath, amssymb}
\usepackage[margin=1in]{geometry}
\setlength{\parskip}{6pt}
\title{QUBO Formulation: Financial Fraud Feature Selection}
\date{}
\begin{document}
\maketitle

\subsection*{Financial Fraud Feature Selection}

\paragraph{Problem Statement.}
Let $N$ denote the number of candidate features. For each feature $i \in \{0, \dots, N-1\}$, let $v_i \in \mathbb{R}$ represent its predictive value. For each pair of distinct features $i, j \in \{0, \dots, N-1\}$ with $i < j$, let $r_{ij} \in \mathbb{R}$ represent the redundancy penalty between them. The objective is to select a subset of exactly $K$ features to maximize the total predictive value minus the total redundancy penalty among selected features. Formally, we seek to maximize
\begin{equation}
    O(\mathbf{x}) = \sum_{i=0}^{N-1} v_i x_i - \sum_{0 \le i < j \le N-1} r_{ij} x_i x_j
\end{equation}
subject to the constraint that exactly $K$ features are selected.

\paragraph{Decision Variables.}
We define binary decision variables $x_i \in \{0,1\}$ for each candidate feature $i \in \{0, \dots, N-1\}$, where $x_i = 1$ indicates that feature $i$ is included in the selection, and $x_i = 0$ indicates it is excluded.

\paragraph{QUBO Derivation.}
\textbf{Step 1: Objective Conversion.}
The Quadratic Unconstrained Binary Optimization (QUBO) framework minimizes a quadratic Hamiltonian $H(\mathbf{x})$. To convert the maximization problem to minimization, we negate the objective function:
\begin{equation}
    H_{\text{obj}}(\mathbf{x}) = -O(\mathbf{x}) = -\sum_{i=0}^{N-1} v_i x_i + \sum_{0 \le i < j \le N-1} r_{ij} x_i x_j.
\end{equation}

\textbf{Step 2: Constraint Formulation.}
The requirement to select exactly $K$ features is expressed as the equality constraint:
\begin{equation}
    \sum_{i=0}^{N-1} x_i = K.
\end{equation}

\textbf{Step 3: Penalty Term Expansion.}
We enforce the constraint by adding a quadratic penalty term with weight $\lambda > 0$:
\begin{equation}
    P(\mathbf{x}) = \lambda \left( \sum_{i=0}^{N-1} x_i - K \right)^2.
\end{equation}
Expanding the square and utilizing the idempotent property $x_i^2 = x_i$ for binary variables:
\begin{align}
    P(\mathbf{x}) &= \lambda \left( \sum_{i=0}^{N-1} x_i^2 + \sum_{i \neq j} x_i x_j - 2K \sum_{i=0}^{N-1} x_i + K^2 \right) \\
    &= \lambda \left( \sum_{i=0}^{N-1} x_i + 2 \sum_{0 \le i < j \le N-1} x_i x_j - 2K \sum_{i=0}^{N-1} x_i + K^2 \right) \\
    &= \sum_{i=0}^{N-1} \lambda(1 - 2K) x_i + \sum_{0 \le i < j \le N-1} 2\lambda x_i x_j + \lambda K^2.
\end{align}

\textbf{Step 4: Combined Hamiltonian.}
The total QUBO Hamiltonian is the sum of the objective and penalty terms:
\begin{align}
    H(\mathbf{x}) &= H_{\text{obj}}(\mathbf{x}) + P(\mathbf{x}) \\
    &= \sum_{i=0}^{N-1} \left( -v_i + \lambda(1 - 2K) \right) x_i + \sum_{0 \le i < j \le N-1} \left( r_{ij} + 2\lambda \right) x_i x_j + \lambda K^2.
\end{align}

\textbf{Step 5: Q Matrix and Offset.}
Reading off the coefficients from $H(\mathbf{x})$, the symmetric Q matrix entries and constant offset $c$ are given by:
\begin{align}
    Q_{i,i} &= -v_i + \lambda(1 - 2K), \quad \forall i \in \{0, \dots, N-1\}, \\
    Q_{i,j} &= r_{ij} + 2\lambda, \quad \forall 0 \le i < j \le N-1, \\
    Q_{j,i} &= Q_{i,j}, \\
    c &= \lambda K^2.
\end{align}

\paragraph{Penalty Weight Justification.}
The penalty weight $\lambda$ must be chosen sufficiently large to ensure that any energy reduction gained from the objective function by violating the constraint is outweighed by the penalty cost. Let $M = \max_{\mathbf{x} \in \{0,1\}^N} \left| \sum_{i=0}^{N-1} v_i x_i - \sum_{0 \le i < j \le N-1} r_{ij} x_i x_j \right|$ be the maximum absolute value of the objective function over the entire hypercube. Choosing $\lambda > M$ guarantees that for any infeasible assignment, the penalty term $P(\mathbf{x})$ strictly exceeds the maximum possible deviation of the objective term from its minimum, ensuring that infeasible solutions cannot achieve a lower total energy than feasible ones.

\paragraph{Correctness Argument.}
Minimizing $H(\mathbf{x})$ over $\{0,1\}^N$ recovers the optimal solution to the original constrained problem. For any feasible assignment $\mathbf{x}_{\text{feas}}$, the penalty term vanishes, yielding $H(\mathbf{x}_{\text{feas}}) = -O(\mathbf{x}_{\text{feas}})$. For any infeasible assignment $\mathbf{x}_{\text{infeas}}$, the penalty term is strictly positive. With $\lambda > M$, the penalty cost $P(\mathbf{x}_{\text{infeas}})$ exceeds the maximum possible improvement in $-O(\mathbf{x})$ achievable by any infeasible solution relative to the best feasible solution. Consequently, $H(\mathbf{x}_{\text{infeas}}) > H(\mathbf{x}_{\text{feas}})$ for all feasible $\mathbf{x}_{\text{feas}}$, implying that the global minimum of $H(\mathbf{x})$ must occur at a feasible assignment that maximizes the original objective function.

\paragraph{Illustrative Example.}
Consider an instance with $N=3$ candidate features and $K=2$ features to be selected. Using the general formulas derived above with penalty weight $\lambda = 1000$, we instantiate the Hamiltonian. The penalty expansion yields diagonal coefficients $\lambda(1-2K) = 1000(1-4) = -3000$, off-diagonal coefficients $2\lambda = 2000$, and offset $\lambda K^2 = 1000(2^2) = 4000$. The resulting QUBO Hamiltonian is:
\begin{align}
    H(\mathbf{x}) &= (-v_0 - 3000)x_0 + (-v_1 - 3000)x_1 + (-v_2 - 3000)x_2 \nonumber \\
    &\quad + (r_{01} + 2000)x_0 x_1 + (r_{02} + 2000)x_0 x_2 + (r_{12} + 2000)x_1 x_2 + 4000.
\end{align}
The corresponding Q matrix entries are $Q_{0,0} = -v_0 - 3000$, $Q_{0,1} = r_{01} + 2000$, and $c = 4000$, with analogous expressions for indices $1$ and $2$. This formulation encodes the specific instance parameters $N=3$, $K=2$, and $\lambda=1000$ while retaining the predictive values $v_i$ and redundancy penalties $r_{ij}$ as problem-specific parameters.
\end{document}